% TOP_PROBLEM: {"id": 2307040, "title": "Research Problems in Function Theory — Problem 7.40", "category_id": 8, "category": {"id": 8, "name": "analysis", "display_name": "Analysis", "description": "Limits, continuity, calculus, and function theory.", "slug": "analysis", "order_index": 8, "created_at": "2026-07-31T15:26:25.671Z"}, "status": "open"}
% FIRST_POSTED: null
% ENTRY_KIND: statement_only
% Imported from: batch05.tex; UnsolvedMath ID 2307040
\documentclass[11pt]{article}
\usepackage{amsmath,amssymb,mathtools}
\usepackage{fontspec,unicode-math}
\setmainfont{Latin Modern Roman}
\setmathfont{Latin Modern Math}
\usepackage{xurl,hyperref}
\newcommand{\sref}[2]{\hyperlink{source:#1:#2}{[S#2]}}
\newcommand{\cataloguescope}[1]{\paragraph{Mathematical statement.}#1}
\begin{document}
% UnsolvedMath ID 2307040; source catalogue code AMR-022-7040
\setcounter{section}{2307039}
\section{Research Problems in Function Theory — Problem 7.40}\label{2307040}
{\small\sffamily UnsolvedMath ID 2307040 · Analysis}
\subsection{Definitions and mathematical statement}

\paragraph{Shared notation.}$\mathbb N=\{1,2,\ldots\}$, and $\mathbb Z,\mathbb Q,\mathbb R,\mathbb C$ denote the integers, rationals, reals and complex numbers. Any other conventions are specified locally.


Let $\overline B(z,r)=\{w\in\mathbb C:|w-z|\le r\}$. For $N$ pairwise disjoint closed discs of the same positive radius, contained in $\{|z|\le1/2\}$, define
\[\Omega=\mathbb D\setminus\bigcup_{j=1}^N\overline B(z_j,r).\]
Let $h$ be its harmonic-measure function: the continuous solution on $\overline\Omega$ of $\Delta h=0$ in $\Omega$, with boundary values one on $\partial\mathbb D$ and zero on the inner circles, extended by zero over the removed discs. Write $\lfloor x\rfloor$ for the greatest integer at most $x$.

\paragraph{Statement recorded in the supplied source.}
Is there an absolute $\delta\in(0,1/4)$ such that, whenever
\[0<r\le\delta,\qquad N\ge\lfloor1/r\rfloor^{\,2-\delta},\]
at least one index $j$ satisfies
\[\int_0^{2\pi}h(z_j+2re^{i\theta})\,d\theta\le r^{2+\delta}?\] \sref{2307040}{2}
Taking $\delta<1/4$ ensures that all circles being sampled stay inside the unit disc and is harmless for this existential small-radius assertion. Update7.40 records Lewis and Wu's proof and its role in their solutions of Problems1.6 and4.18. \sref{2307040}{2}
\subsection{Short English statement}
When sufficiently many equal small discs are removed, must at least one have an exceptionally small surrounding circular integral of the outer-boundary harmonic measure?
\subsection{Sources}
\begin{description}
\item[\hypertarget{source:2307040:1}{S1}] ULAM.AI and the UnsolvedMath Contributors, \emph{UnsolvedMath: A Curated Collection of Open Mathematics Problems}, record 2307040 (AMR-022-7040). CC BY4.0. \url{https://huggingface.co/datasets/ulamai/UnsolvedMath}.
\item[\hypertarget{source:2307040:2}{S2}] W. K. Hayman and E. F. Lingham, \emph{Research Problems in Function Theory} (2018), Chapter7, Problem7.40 and its complete update. \url{https://arxiv.org/abs/1809.07200}.
\end{description}
\label{end:2307040}
\end{document}
